\section{Core Concepts and \acr{LLVM} Syntax}

\acr{LLVM} \acr{IR} is a low-level, strongly-typed intermediate language that resembles assembly but provides platform independence and optimization opportunities. This section introduces the fundamental concepts and syntax elements that form the foundation of \acr{LLVM} \acr{IR} programming.

\subsection{Static Single Assignment (\acr{SSA}) Form}

\acr{SSA} is a property where each variable is assigned exactly once and every use refers to a specific definition. This constraint makes data flow explicit and enables powerful compiler optimizations by simplifying the analysis of how values propagate through the program.

\textbf{Non-\acr{SSA} code:}
\begin{lstlisting}[style=ir]
x = 1
x = x + 2  // x is assigned twice (not SSA)
\end{lstlisting}

\textbf{\acr{SSA} code:}
\begin{lstlisting}[style=ir]
%x1 = 1
%x2 = add i32 %x1, 2  // Each value has unique name
\end{lstlisting}

The benefits of \acr{SSA} form include:
\begin{itemize}[noitemsep]
  \item \textbf{Constant propagation}: The value of \texttt{\%x2} is provably 3, enabling compile-time evaluation
  \item \textbf{Dead code elimination}: Unused definitions can be safely removed without complex liveness analysis
  \item \textbf{Register allocation}: Clear lifetime of each value simplifies mapping virtual registers to physical registers
  \item \textbf{Simplified dependency tracking}: Each definition has a unique name, making def-use chains trivial to construct
\end{itemize}

\subsection{Modules and Target Specification}

\acr{LLVM} \acr{IR} is organized hierarchically, with the \textbf{module} serving as the top-level container. A module corresponds to a single compilation unit and contains functions, global variables, type definitions, and metadata.

Every module must specify its target architecture through two declarations:

\begin{lstlisting}[style=ir]
target triple = "x86_64-unknown-linux-gnu"
target datalayout = "e-m:e-p270:32:32-i64:64-f80:128-n8:16:32:64-S128"
\end{lstlisting}

The \textbf{target triple} follows the format \texttt{architecture-vendor-os-environment} and identifies the platform:

\begin{itemize}[noitemsep]
  \item \texttt{\acr{x86}\_64-apple-darwin} --- macOS on Intel
  \item \texttt{aarch64-apple-darwin} --- macOS on Apple Silicon
  \item \texttt{aarch64-unknown-linux-gnu} --- Linux on \acr{ARM} (\acr{AWS} Graviton, Raspberry Pi)
  \item \texttt{nvptx64-nvidia-cuda} --- NVIDIA \acr{GPU}s
  \item \texttt{amdgcn-amd-amdhsa} --- \acr{AMD} \acr{GPU}s (\acr{ROCm})
  \item \texttt{wasm32-unknown-unknown} --- \acr{WASM}
  \item \texttt{spirv64-unknown-unknown} --- \acr{SPIR-V} for Vulkan/\acr{OpenCL}
\end{itemize}

The \textbf{data layout string} is a compact specification that encodes essential low-level architectural properties required for correct code generation~\cite{llvmlangref}. This string informs the \acr{LLVM} optimizer and backend about target-specific characteristics such as endianness, pointer sizes, type alignments, and memory layout conventions~\cite{lopesauler2014}. Without accurate data layout information, the compiler cannot generate correct memory access patterns, properly align data structures, or make sound optimization decisions that depend on type sizes and alignment constraints~\cite{llvmessentials}.

\subsubsection*{Data Layout String Format}

The data layout specification consists of multiple components separated by the minus sign (\texttt{-}), where each component begins with a letter identifying the specification type, followed by parameters~\cite{llvmlangref}. A typical \acr{x86}-64 data layout string is:

\begin{lstlisting}[style=ir]
target datalayout = "e-m:e-p270:32:32-p271:32:32-p272:64:64-i64:64-f80:128-n8:16:32:64-S128"
\end{lstlisting}

\subsubsection*{Component Breakdown}

\textbf{Endianness (\texttt{e} or \texttt{E}):} The first character specifies byte ordering~\cite{llvmlangref}:
\begin{itemize}[noitemsep]
  \item \texttt{e} --- Little-endian: least significant byte at lowest memory address (Intel/\acr{AMD} \acr{x86}, \acr{ARM} in little-endian mode)
  \item \texttt{E} --- Big-endian: most significant byte at lowest memory address (legacy PowerPC, SPARC, some network protocols)
\end{itemize}

Endianness affects how multi-byte values are stored in memory. For example, the 32-bit integer \texttt{0x12345678} is stored as:
\begin{itemize}[noitemsep]
  \item Little-endian: \texttt{78 56 34 12} (bytes at increasing addresses)
  \item Big-endian: \texttt{12 34 56 78} (bytes at increasing addresses)
\end{itemize}

\textbf{Name Mangling (\texttt{m:<style>}):} Specifies how symbol names are mangled for the linker~\cite{llvmlangref}:
\begin{itemize}[noitemsep]
  \item \texttt{m:e} --- ELF mangling (Linux, BSD, most Unix systems)
  \item \texttt{m:o} --- Mach-O mangling (macOS, iOS)
  \item \texttt{m:m} --- Mips mangling
  \item \texttt{m:w} --- Windows COFF mangling
  \item \texttt{m:x} --- Windows COFF \acr{x86} mangling (different from \texttt{m:w} for compatibility)
\end{itemize}

Symbol mangling affects how function and global variable names are encoded in object files, determining linker compatibility and calling convention details.

\textbf{Pointer Layout (\texttt{p[n]:<size>:<abi>:<pref>}):} Specifies pointer characteristics for address space \texttt{n}~\cite{llvmlangref}. If \texttt{n} is omitted, address space 0 (default) is assumed:
\begin{itemize}[noitemsep]
  \item \texttt{<size>} --- Pointer size in bits
  \item \texttt{<abi>} --- \acr{ABI}-required alignment in bits (minimum guaranteed by calling conventions)
  \item \texttt{<pref>} --- Preferred alignment in bits for performance (optional; if omitted, equals \texttt{<abi>})
\end{itemize}

Examples from the \acr{x86}-64 data layout:
\begin{itemize}[noitemsep]
  \item \texttt{p270:32:32} --- Address space 270: 32-bit pointers with 32-bit alignment (used for \texttt{\_\_ptr32} Microsoft extension)
  \item \texttt{p271:32:32} --- Address space 271: 32-bit zero-extended pointers (used for \texttt{\_\_uptr})
  \item \texttt{p272:64:64} --- Address space 272: 64-bit pointers with 64-bit alignment (used for \texttt{\_\_ptr64})
  \item Default (\texttt{p:64:64:64}) --- Address space 0: 64-bit pointers with 64-bit alignment (standard pointers)
\end{itemize}

Multiple address spaces enable heterogeneous programming where different pointer types have different sizes (e.g., \acr{GPU} programming with separate host and device address spaces)~\cite{lopesauler2014}.

\textbf{Integer Alignment (\texttt{i<size>:<abi>:<pref>}):} Specifies alignment requirements for integer types~\cite{llvmlangref}:

\begin{lstlisting}[style=ir]
i64:64      ; i64 has 64-bit ABI and preferred alignment
i32:32      ; i32 has 32-bit alignment
i8:8        ; i8 has 8-bit alignment
\end{lstlisting}

Correct alignment is critical for performance and correctness. Misaligned memory access can cause:
\begin{itemize}[noitemsep]
  \item Performance degradation (multiple memory transactions on \acr{x86})
  \item Hardware exceptions or crashes (strict-alignment architectures like older \acr{ARM}, SPARC)
  \item Incorrect atomic operations (atomics require natural alignment)
\end{itemize}

\textbf{Floating-Point Alignment (\texttt{f<size>:<abi>:<pref>}):} Specifies alignment for floating-point types~\cite{llvmlangref}:

\begin{lstlisting}[style=ir]
f80:128     ; x86 80-bit extended precision (stored in 128-bit slots)
f64:64      ; double (64-bit float)
f32:32      ; float (32-bit float)
\end{lstlisting}

The \texttt{f80:128} entry indicates that while \acr{x86} extended precision floats are 80 bits, they are stored in 128-bit-aligned memory slots for efficiency.

\textbf{Vector Alignment (\texttt{v<size>:<abi>:<pref>}):} Specifies alignment for vector types used in \acr{SIMD} operations~\cite{llvmlangref}:

\begin{lstlisting}[style=ir]
v128:128    ; 128-bit vectors (SSE on x86, NEON on ARM)
v256:256    ; 256-bit vectors (AVX)
v512:512    ; 512-bit vectors (AVX-512)
\end{lstlisting}

Proper vector alignment is mandatory for \acr{SIMD} instruction efficiency. Misaligned vector loads/stores can cause significant performance penalties or instruction faults.

\textbf{Aggregate Alignment (\texttt{a:<abi>:<pref>}):} Specifies alignment for aggregate types (structs, arrays)~\cite{llvmlangref}:

\begin{lstlisting}[style=ir]
a:0:64      ; Aggregates have 0-bit ABI alignment, 64-bit preferred alignment
\end{lstlisting}

This allows the compiler to optimize struct layout while maintaining \acr{ABI} compatibility.

\textbf{Native Integer Widths (\texttt{n<size1>:<size2>:...}):} Specifies integer widths that the target CPU can efficiently process in a single instruction~\cite{llvmlangref}:

\begin{lstlisting}[style=ir]
n8:16:32:64 ; x86-64 efficiently supports 8, 16, 32, and 64-bit integers
n32:64      ; PowerPC 64 efficiently supports 32 and 64-bit integers
\end{lstlisting}

The optimizer uses this information for type legalization decisions. Operations on \textbf{\textit{non-native widths}} (e.g., \texttt{i7}, \texttt{i128} on \acr{x86}-64) are decomposed into multiple native-width operations.

\textbf{Stack Alignment (\texttt{S<size>}):} Specifies the natural stack alignment in bits~\cite{llvmlangref}:

\begin{lstlisting}[style=ir]
S128        ; Stack maintains 128-bit (16-byte) alignment
\end{lstlisting}

Stack alignment affects:
\begin{itemize}[noitemsep]
  \item Function prologues/epilogues (stack frame setup)
  \item \acr{SIMD} variable allocation (local vectors must be aligned)
  \item Calling conventions (some \acr{ABI}s require specific stack alignment)
\end{itemize}

On \acr{x86}-64, 16-byte stack alignment enables efficient \acr{SSE} operations on stack-allocated vectors~\cite{lopesauler2014}.

The data layout string is essential for three fundamental reasons~\cite{llvmessentials}: \textbf{(1) Correctness}---wrong alignment can cause crashes on strict-alignment architectures or incorrect atomic operations; \textbf{(2) Performance}---proper alignment enables efficient memory access, \acr{SIMD} vectorization, and cache utilization; and \textbf{(3) \acr{ABI} compliance}---data layout must match the platform \acr{ABI} to ensure interoperability with system libraries and other compiled code. Modern \acr{LLVM} automatically sets the data layout when you specify the target triple, but understanding the format is crucial for cross-compilation, \acr{GPU} programming, and custom target development~\cite{lopesauler2014}.

\subsection{Functions}

Functions are the primary units of executable code in \acr{LLVM} \acr{IR}. A function declaration consists of:

\begin{lstlisting}[style=ir]
define <return_type> @function_name(<param_type> %param1, ...) {
  ; function body
}
\end{lstlisting}

\textbf{Linkage types} control visibility and linking behavior:

\begin{itemize}[noitemsep]
  \item \texttt{define} --- Default linkage (internal to module unless exported)
  \item \texttt{define external} --- Exported symbol, visible to other modules
  \item \texttt{define internal} --- Private to the module (similar to C \texttt{static})
  \item \texttt{define weak} --- Can be overridden by other definitions
  \item \texttt{declare} --- External function (declaration only, no body)
\end{itemize}

\textbf{Function attributes} modify compilation behavior:

\begin{lstlisting}[style=ir]
define i32 @pure_function(i32 %x) nounwind readnone {
  %result = mul i32 %x, %x
  ret i32 %result
}
\end{lstlisting}

Common attributes:
\begin{itemize}[noitemsep]
  \item \texttt{nounwind} --- Function never throws exceptions
  \item \texttt{readnone} --- Pure function, no memory access
  \item \texttt{readonly} --- Function only reads memory
  \item \texttt{alwaysinline} --- Force inlining
  \item \texttt{noinline} --- Prevent inlining
\end{itemize}

\subsection{Basic Blocks and Labels}

\acr{LLVM} organizes function bodies into \textbf{basic blocks}---straight-line sequences of instructions with no branches except at the end. This structure is fundamental to control flow analysis and optimization.

\begin{lstlisting}[style=ir]
define i32 @max(i32 %a, i32 %b) {
entry:
  %cmp = icmp sgt i32 %a, %b  ; Compare: a > b?
  br i1 %cmp, label %if.then, label %if.else

if.then:
  ret i32 %a

if.else:
  ret i32 %b
}
\end{lstlisting}

Each basic block has three essential properties:

\begin{enumerate}[noitemsep]
  \item \textbf{Unique label}: Every block has a name (\texttt{entry}, \texttt{if.then}, \texttt{if.else}) that serves as a branch target
  \item \textbf{Sequential instructions}: No control flow transfers occur in the middle of a block
  \item \textbf{Terminator instruction}: Every block must end with exactly one terminator that transfers control
\end{enumerate}

\textbf{Terminator instructions:}

\begin{itemize}[noitemsep]
  \item \texttt{ret <type> <value>} --- Return from function
  \item \texttt{br label \%dest} --- Unconditional jump
  \item \texttt{br i1 \%cond, label \%true, label \%false} --- Conditional branch
  \item \texttt{switch} --- Multi-way branch (jump table)
  \item \texttt{unreachable} --- Marks dead code paths
\end{itemize}

The first basic block in every function is conventionally named \texttt{entry} and serves as the function's entry point.

\subsection{Global and Local Variables}

\acr{LLVM} distinguishes between global variables (module-level storage) and local variables (function-level virtual registers).

\textbf{Global variables} are declared outside functions and persist for the program's lifetime:

\begin{lstlisting}[style=ir]
@global_counter = global i32 0
@constant_pi = constant float 3.14159
@external_array = external global [100 x i32]
\end{lstlisting}

Global variables use the \texttt{@} prefix and have linkage specifiers:
\begin{itemize}[noitemsep]
  \item \texttt{global} --- Mutable global variable
  \item \texttt{constant} --- Immutable global constant
  \item \texttt{external} --- Declared elsewhere (no initializer)
  \item \texttt{internal} --- Private to the module
\end{itemize}

\textbf{Local variables} (virtual registers) use the \texttt{\%} prefix and exist only within a function:

\begin{lstlisting}[style=ir]
define i32 @compute(i32 %x) {
entry:
  %temp1 = add i32 %x, 5
  %temp2 = mul i32 %temp1, 2
  ret i32 %temp2
}
\end{lstlisting}

In \acr{SSA} form, each local variable is assigned exactly once. For mutable variables (stack allocations), \acr{LLVM} uses memory operations:

\begin{lstlisting}[style=ir]
define void @mutable_example() {
entry:
  %ptr = alloca i32              ; Allocate stack space
  store i32 42, ptr %ptr         ; Write to memory
  %val = load i32, ptr %ptr      ; Read from memory
  ret void
}
\end{lstlisting}

\subsection{Type System}

\acr{LLVM} \acr{IR} is strongly typed. Every value has an explicit type, and type mismatches are compile-time errors. The type system includes primitive types, aggregate types, and pointer types.

\textbf{Primitive types:}

\begin{itemize}[noitemsep]
  \item \textbf{Integer types}: \texttt{i1} (boolean), \texttt{i8}, \texttt{i16}, \texttt{i32}, \texttt{i64}, \texttt{i128} (arbitrary-width integers like \texttt{i7} or \texttt{i256} are also valid)
  \item \textbf{Floating-point types}: \texttt{half} (16-bit), \texttt{float} (32-bit), \texttt{double} (64-bit), \texttt{fp128} (quad-precision)
  \item \textbf{Void type}: \texttt{void} (used only for function return types, not for values)
\end{itemize}

\textbf{Aggregate types:}

\begin{lstlisting}[style=ir]
; Array: [length x element_type]
[10 x i32]                    ; Array of 10 i32 values
[3 x float]                   ; Array of 3 floats

; Structure: { type1, type2, ... }
{ i32, float, ptr }           ; Unpacked struct
<{ i8, i32 }>                ; Packed struct (no padding)

; Vector: <length x element_type>
<4 x float>                   ; SIMD vector (SSE, NEON)
<16 x i8>                     ; 128-bit packed integer vector
\end{lstlisting}

\textbf{Pointer type}: \texttt{ptr} (opaque pointer, introduced in \acr{LLVM} 15):

\begin{lstlisting}[style=ir]
%heap_ptr = call ptr @malloc(i64 64)
store i32 42, ptr %heap_ptr
%value = load i32, ptr %heap_ptr
\end{lstlisting}

Older \acr{LLVM} versions used typed pointers (\texttt{i32*}, \texttt{float*}), but modern \acr{LLVM} uses opaque \texttt{ptr} to simplify the type system and improve optimization.

\textbf{Function types} specify signatures:

\begin{lstlisting}[style=ir]
; Function type: return_type (param_types)
i32 (i32, i32)                ; Takes two i32, returns i32
void (ptr, i64)               ; Takes pointer and i64, returns void
\end{lstlisting}

\subsection{Instructions and Opcodes}

\acr{LLVM} instructions are the fundamental operations that manipulate values. Every instruction produces a result (except terminators and side-effecting operations like \texttt{store}).

\textbf{Arithmetic instructions:}

\begin{lstlisting}[style=ir]
%sum = add i32 %a, %b         ; Integer addition
%diff = sub i32 %a, %b        ; Integer subtraction
%prod = mul i32 %a, %b        ; Integer multiplication
%quot = sdiv i32 %a, %b       ; Signed division
%rem = srem i32 %a, %b        ; Signed remainder

%fsum = fadd float %x, %y     ; Floating-point addition
%fdiff = fsub float %x, %y    ; Floating-point subtraction
%fprod = fmul float %x, %y    ; Floating-point multiplication
%fquot = fdiv float %x, %y    ; Floating-point division
\end{lstlisting}

\textbf{Bitwise operations:}

\begin{lstlisting}[style=ir]
%and_result = and i32 %a, %b  ; Bitwise AND
%or_result = or i32 %a, %b    ; Bitwise OR
%xor_result = xor i32 %a, %b  ; Bitwise XOR
%shift_left = shl i32 %a, 2   ; Shift left by 2 bits
%shift_right = lshr i32 %a, 2 ; Logical shift right
%arith_shift = ashr i32 %a, 2 ; Arithmetic shift right (sign-extend)
\end{lstlisting}

\textbf{Comparison instructions:}

\begin{lstlisting}[style=ir]
; Integer comparisons (icmp)
%eq = icmp eq i32 %a, %b      ; Equal
%ne = icmp ne i32 %a, %b      ; Not equal
%gt = icmp sgt i32 %a, %b     ; Signed greater than
%lt = icmp slt i32 %a, %b     ; Signed less than
%uge = icmp uge i32 %a, %b    ; Unsigned greater or equal

; Floating-point comparisons (fcmp)
%feq = fcmp oeq float %x, %y  ; Ordered equal
%fgt = fcmp ogt float %x, %y  ; Ordered greater than
%fne = fcmp une float %x, %y  ; Unordered or not equal
\end{lstlisting}

\textbf{Memory operations:}

\begin{lstlisting}[style=ir]
%ptr = alloca i32             ; Stack allocation
store i32 42, ptr %ptr        ; Write to memory
%val = load i32, ptr %ptr     ; Read from memory

; Pointer arithmetic with getelementptr (GEP)
%array_ptr = getelementptr [10 x i32], ptr %base, i64 0, i64 3
; Access element 3 of array at %base
\end{lstlisting}

\textbf{Type conversion instructions:}

\begin{lstlisting}[style=ir]
%ext = zext i32 %a to i64     ; Zero-extend to wider type
%trunc = trunc i64 %b to i32  ; Truncate to narrower type
%fp_to_int = fptosi float %x to i32  ; Float to signed int
%int_to_fp = sitofp i32 %y to float  ; Signed int to float
%bitcast = bitcast i32 %z to float   ; Reinterpret bits
\end{lstlisting}

\textbf{Function calls:}

\begin{lstlisting}[style=ir]
%result = call i32 @compute(i32 %x, i32 %y)
call void @print_value(i32 %result)
\end{lstlisting}

\textbf{Vector operations (\acr{SIMD}):}

\begin{lstlisting}[style=ir]
%vec_a = <4 x float> <float 1.0, float 2.0, float 3.0, float 4.0>
%vec_b = <4 x float> <float 5.0, float 6.0, float 7.0, float 8.0>
%vec_sum = fadd <4 x float> %vec_a, %vec_b
\end{lstlisting}

\subsection{The \acr{PHI} Instruction}

When control flow merges from multiple basic blocks, \acr{SSA} form requires a mechanism to select the appropriate value based on which predecessor block was executed. The \texttt{phi} instruction serves this purpose.

\begin{lstlisting}[style=ir]
define i32 @abs(i32 %x) {
entry:
  %is_neg = icmp slt i32 %x, 0
  br i1 %is_neg, label %negate, label %done

negate:
  %neg = sub i32 0, %x
  br label %done

done:
  %result = phi i32 [ %x, %entry ], [ %neg, %negate ]
  ret i32 %result
}
\end{lstlisting}

The \texttt{phi} instruction syntax is:

\begin{lstlisting}[style=ir]
%result = phi <type> [ <value1>, %block1 ], [ <value2>, %block2 ], ...
\end{lstlisting}

Semantics: "If control flow arrived from \texttt{\%block1}, use \texttt{<value1>}; if from \texttt{\%block2}, use \texttt{<value2>}."

\textbf{Why \texttt{phi} is necessary:}

Without \texttt{phi}, we cannot maintain \acr{SSA} form when values merge from different control paths. Consider this invalid attempt:

\begin{lstlisting}[style=ir]
; INVALID: Cannot assign to %result in multiple blocks
entry:
  %result = ... ; First assignment
  br ...

other_block:
  %result = ... ; ERROR: violates SSA (redefinition)
\end{lstlisting}

The \texttt{phi} instruction solves this by creating a single assignment point that selects from multiple incoming values.

\textbf{Multiple \texttt{phi} nodes:}

A block can have multiple \texttt{phi} instructions. All \texttt{phi} instructions in a block are conceptually executed simultaneously before any other instructions:

\begin{lstlisting}[style=ir]
loop:
  %i = phi i32 [ 0, %entry ], [ %i_next, %loop ]
  %sum = phi i32 [ 0, %entry ], [ %sum_next, %loop ]
  %i_next = add i32 %i, 1
  %sum_next = add i32 %sum, %i
  %cond = icmp slt i32 %i_next, 10
  br i1 %cond, label %loop, label %exit
\end{lstlisting}

This example implements a loop that sums integers from 0 to 9, demonstrating how \texttt{phi} enables iterative control flow in \acr{SSA} form.
